\itemhead{习题2.1 · 第11题（教材p.44）}{XJ-B1-C02-HW-P044-Y1102}
\label{XJ-B1-C02-HW-P044-Y1102}
{\footnotesize 湘教版教材p.44 / PDF第51页 / 习题2.1 / 第11题}
\field{原题}{某商品计划提价两次，有甲、乙、丙三种方案，其中 $p>q>0$。经两次提价后，哪种方案提价的幅度大？为什么？
\begin{center}\begin{tabular}{c|c|c}\hline
方案&第一次提价&第二次提价\\\hline
甲&$p\%$&$q\%$\\
乙&$q\%$&$p\%$\\
丙&$\dfrac{p+q}{2}\%$&$\dfrac{p+q}{2}\%$\\\hline
\end{tabular}\end{center}}
\field{参考答案}{甲、乙幅度相同；丙最大。令 $P=p/100,Q=q/100$，丙与甲的价格倍率差为 $[(P+Q)/2]^2-PQ=(P-Q)^2/4>0$。}
{\small \field{关联知识}{I03 基本不等式、等号与最值建模（教材帮p.52、53、54 / PDF61、62、63）}
\field{方法与易错点}{列正数条件与固定量，构造基本不等式，解等号条件并代回实际限制。 易错：把一个下界误当最小值；多个等号条件不能同时实现；漏单位和整数限制。}
\field{答案依据}{答案册相关题组：PDF第12、13页（册内页11、12）；按当前题库原题逐项重算，不按旧版题号机械对套。}}
